\subsubsection{Momentos de inércia}
\label{sec:dinAcoplada_momentos_inercia}

A matriz de inércia generalizada $M(\vec q)$ é introduzida a partir da forma quadrática da energia cinética, conforme \cite{spong2006robot}:

\begin{equation}
T = \tfrac{1}{2}\,\dot{\vec q}^{\,T} M(\vec q)\,\dot{\vec q}.
\end{equation}

Adota-se o vetor de coordenadas generalizadas \cite{bloesch2013}:

\begin{equation}
\vec q = [\vec q_B^{\,T}\;\;\vec\theta^{\,T}]^{\,T},
\end{equation}

no qual $\vec q_B$ descreve a atitude da base e $\vec\theta$ reúne as juntas do \gls{mr}.
Com essa escolha, a matriz $M(\vec q)$ assume a forma particionada usual para sistemas de base flutuante, conforme \cite{featherstone2008}:

\begin{equation}
M(\vec q) =
\begin{bmatrix}
M_{bb} & M_{bm} \\[4pt]
M_{mb} & M_{mm}
\end{bmatrix},
\label{eq:M_blocos}
\end{equation}

sendo que $M_{mb}=M_{bm}^{T}$.

A estrutura em blocos segue da forma como as velocidades espaciais dos elos são escritas a partir dos jacobianos correspondentes \cite{featherstone2008}, $\vec V_i = J_i(\vec q)\,\dot{\vec q}$, onde a energia cinética de cada elo, escrita em termos das velocidades linear e angular do centro de massa \cite{sciavicco_siciliano2010}, pode ser reescrita na forma compacta:


\begin{equation}
T_i
=
\tfrac{1}{2}\,
\vec V_i^{\mathsf T} M_{i,\text{spatial}}\,\vec V_i,
\end{equation}

o que conduz à representação

\begin{equation}
M(\vec q)
=
\sum_i
J_i^{\mathsf T}(\vec q)\,
M_{i,\text{spatial}}\,
J_i(\vec q).
\end{equation}

Em sistemas base--\gls{mr} em microgravidade, essa forma coincide com a matriz de inércia generalizada obtida pela formulação do jacobiano generalizado.
Os blocos de $M(\vec q)$ possuem interpretação direta: $M_{bb}$ reúne as inércias equivalentes da base, $M_{mm}$ contém as contribuições do manipulador, e $M_{bm}$, $M_{mb}$ representam os termos de acoplamento.
Separando \eqref{eq:dinamica_geral} em blocos, obtém-se

\begin{align}
M_{bb}\,\ddot{\vec q}_B
+ M_{bm}\,\ddot{\boldsymbol\theta}
+ C_{bb}\,\dot{\vec q}_B
+ C_{bm}\,\dot{\boldsymbol\theta}
&=
\vec\tau_B,
\\[4pt]
M_{mb}\,\ddot{\vec q}_B
+ M_{mm}\,\ddot{\boldsymbol\theta}
+ C_{mb}\,\dot{\vec q}_B
+ C_{mm}\,\dot{\boldsymbol\theta}
&=
\vec\tau_m,
\end{align}

onde a primeira equação descreve a atitude da base sob as reações do manipulador e a segunda descreve a dinâmica das juntas em base flutuante.
Essa estrutura blocada é análoga àquela obtida por \citeauthor{Wilde2018} para sistemas espaçonave--\gls{mr} em microgravidade, em que a matriz de inércia é escrita na forma

\begin{equation}
\begin{bmatrix}
H_0 & H_{0m}\\
H_{0m}^{\mathsf T} & H_m
\end{bmatrix},
\end{equation}

enquanto que as equações de movimento são separadas em componentes da base e do \gls{mr} \cite[Eq.~(39)]{Wilde2018}.

